% Two workloads on one axis: the node's real rate, and the benchmark.
\begin{tikzpicture}[font=\sffamily\scriptsize]

% ---------------- the node workload
\node[font=\sffamily\bfseries\scriptsize, anchor=west] at (0mm,46mm)
  {The node workload: one block every 256 ms at the 1 kHz rate of chapter 6};
\draw[tspan] (4mm,41mm) -- node[tlabel]{compute: not measured} (13mm,41mm);
\lane{34}{sampling}
\slice{34}{0}{4}{taskisr}{fill}
\slice{34}{80}{4}{taskisr}{fill}
\lane{26}{pipeline}
\slice{26}{4}{9}{taskrun}{run}
\slice{26}{13}{67}{taskblock}{idle, most of the period}
\slice{26}{84}{9}{taskrun}{run}
\draw[tspan] (0mm,22mm) -- node[tlabel]{256 ms block period} (80mm,22mm);
\draw[deadline] (80mm,18mm) -- (80mm,40mm);
\node[font=\tiny, anchor=south west] at (81mm,40mm) {next block due};
\draw[taxis] (0mm,18mm) -- (116mm,18mm) node[right, font=\tiny] {time};

% ---------------- the benchmark workload
\node[font=\sffamily\bfseries\scriptsize, anchor=west] at (0mm,8mm)
  {The benchmark workload: the same transform back to back, nothing else running};
\lane{-2}{pipeline}
\foreach \i in {0,1,2,3,4,5}{
  \slice{-2}{\i*19}{18}{taskrun}{transform}}
\draw[tspan] (0mm,-6mm) -- node[tlabel]{cycles per transform, median of 32} (18mm,-6mm);
\draw[taxis] (0mm,-10mm) -- (116mm,-10mm) node[right, font=\tiny] {time};

% ---------------- notes, clear of both traces
\node[umlnote, text width=52mm, anchor=north west] at (0mm,-16mm)
  {The upper trace is what the node does and it proves the budget. The lower
   trace is what the bench does and it produces the number. Reporting only one
   of the two is how a reader ends up either over-engineering or learning
   nothing about the part.};
\node[umlnote, text width=52mm, anchor=north west] at (62mm,-16mm)
  {The dashed line is the next block arriving. If the run slice ever reaches it
   the node is late, and at this rate on this part that would mean something
   else is wrong.};
\end{tikzpicture}
